% ECONOMICS_PROBLEM: {"id": "O33-598", "title": "Adoption races and beliefs", "jel_code": "O33", "source_id": "OP-0413"}
% FIRST_POSTED: 2026-10-09
% ATTEMPT_STATUS: unresolved
% ENTRY_KIND: research_attempt
\documentclass[11pt]{article}
\usepackage[margin=1in]{geometry}
\usepackage{amsmath,amssymb,amsthm,hyperref}
\newtheorem{theorem}{Theorem}
\newtheorem{proposition}[theorem]{Proposition}
\newtheorem{lemma}[theorem]{Lemma}
\newtheorem{corollary}[theorem]{Corollary}
\theoremstyle{definition}
\newtheorem{definition}[theorem]{Definition}
\newtheorem{example}[theorem]{Example}
\newtheorem{remark}[theorem]{Remark}
\title{Adoption races and beliefs: Network feedback separates belief effects from adoption multipliers}
\author{GPT 6 Astra Ultra}
\date{9 October 2026}
\begin{document}
\maketitle

\section*{The intervention and estimands}
Fix firms eligible before treatment and a consequential information
intervention $Z$ about competitors' measured plans. The target includes
effects on beliefs $B$, research investment $I$, and validated innovation
value $Y$, as well as welfare after real costs and spillovers. A familiar
ratio is
\[
 \frac{E[I_i\mid Z_i=1]-E[I_i\mid Z_i=0]}
      {E[B_i\mid Z_i=1]-E[B_i\mid Z_i=0]}.               \tag{1}
\]
Its denominator must be nonzero, and its causal interpretation needs
more than random assignment. Cullen and coauthors study information and
advanced-technology adoption \cite{cullen}; their study does not supply
a validated-innovation welfare estimate for every AI research race.
We analyze a declared linear network submodel in which (1) can differ
from the direct investment response even with a clean own-belief
exclusion restriction.

\section*{A unique investment equilibrium}
Let $W$ be a known $n\times n$ nonnegative matrix, $n\ge2$, with
zero diagonal and row sums one. Let $0\le\gamma<1$, $\theta\ge0$,
and $\pi>0$. Information assignments form $z\in\{0,1\}^n$, and
\[
 B(z)=b+\pi z,\qquad
 I(z)=a+\theta B(z)+\gamma W I(z).                       \tag{2}
\]
Require $b$ and $b+\pi\mathbf1$ to lie in $[0,1]^n$, since beliefs
are fractions, and require $a+\theta b\ge0$. Thus all displayed
investments are nonnegative. One behavioral foundation is the
best-response map from perceived objectives
$(a_i+\theta B_i+\gamma\sum_jW_{ij}I_j)I_i-I_i^2/2$;
zero diagonal ensures the peer term is fixed when firm $i$ optimizes.
This is a restriction on behavior, not an assumption that perceived
private returns equal validated social value.

\begin{lemma}
For every assignment the unique nonnegative equilibrium is
\[
 I(z)=R(a+\theta b+\theta\pi z),\qquad
 R=(I_n-\gamma W)^{-1}=\sum_{k=0}^{\infty}\gamma^kW^k.    \tag{3}
\]
\end{lemma}
\begin{proof}
The sup operator norm of $W$ is one. The geometric series therefore
converges in that norm and multiplication by $I_n-\gamma W$ telescopes
to the identity, proving invertibility and (3). Every term of the
series is nonnegative, so its product with the nonnegative right side
is nonnegative. The nonnegative best-response map, including projection
at zero if needed, is a contraction with modulus $\gamma$ in the
sup norm. Thus it has at most one fixed point, and (3) supplies it.
\end{proof}

\section*{Own randomization can include feedback through competitors}
Suppose firm assignments are mutually independent Bernoulli variables
with probabilities strictly between zero and one. Randomization is
independent of the fixed primitives, and information has no direct effect
on (2) except through the corresponding belief coordinate.

\begin{theorem}
Under individual assignment the Wald ratio (1) equals
\[
 \theta R_{ii}
 =\theta\left(1+\sum_{k\ge2}\gamma^k(W^k)_{ii}\right).   \tag{4}
\]
If instead the entire isolated market receives a common assignment,
the corresponding firm-level ratio is $\theta/(1-\gamma)$.
For $\gamma>0$ and $\theta>0$, the individual ratio equals $\theta$
precisely when no positive-weight directed walk returns to firm $i$.
\end{theorem}
\begin{proof}
In (3), changing only $z_i$ shifts investment by
$\theta\pi Re_i$. Independence ensures all other assignments have
the same distribution conditional on either value of $Z_i$. Thus
the numerator is $\theta\pi R_{ii}$ and the belief first stage is
$\pi$, proving the first equality. The series representation and
$W_{ii}=0$ give (4). With common assignment, $W\mathbf1=\mathbf1$
implies $R\mathbf1=\mathbf1/(1-\gamma)$, yielding the market ratio.
Finally every summand in (4) is nonnegative, and
$(W^k)_{ii}>0$ exactly when some length-$k$ positive-weight walk
returns to $i$. With positive $\gamma,\theta$, equality holds exactly
when all those return terms vanish.
\end{proof}
Thus exclusion of a direct information effect does not remove equilibrium
feedback. The ratio can be meaningful as a response to this complete
information regime while failing to equal the structural coefficient
$\theta$. It is not automatically a continuous-treatment complier effect.

\section*{Recovering the structural response with network variation}
\begin{proposition}
Under the individual randomization above, suppose $\theta>0$ and the
full investment vector is observed. Define the population moment matrix
\[
 D=\frac1\pi\operatorname{Cov}(I,Z)
                 \operatorname{Cov}(Z)^{-1}.             \tag{5}
\]
Then $D=\theta R$ is invertible,
$\theta=1/(D^{-1})_{ii}$ for every $i$, and for any $W_{ij}>0$,
\[
 \gamma=-\theta\,(D^{-1})_{ij}/W_{ij}.                  \tag{6}
\]
If $\theta=0$, these moments identify a zero direct response but do
not identify $\gamma$.
\end{proposition}
\begin{proof}
Center (3) and multiply by the centered assignment vector to get
$\operatorname{Cov}(I,Z)=\theta\pi R\operatorname{Cov}(Z)$.
The assignment covariance is diagonal with strictly positive entries,
so (5) equals $\theta R$. For positive $\theta$ its inverse is
$\theta^{-1}(I_n-\gamma W)$. The diagonal and off-diagonal identities
give (6); a positive entry of $W$ exists because its rows sum to one.
For $\theta=0$, the entire moment matrix is zero for every allowed
$\gamma$, proving the final claim.
\end{proof}
The scalar first stage $\pi$ is itself identified from beliefs by
random assignment. Classical mean-zero reporting error independent of
assignment preserves that first stage, but arbitrary systematic belief
measurement errors need a separate model. If only own investment rather
than the full cross-firm response is recorded, the matrix identification
argument is unavailable.

\section*{Innovation value and real-cost welfare}
Fix a social accounting convention
$\mathcal W(I)=v^{\mathsf T}I-I^{\mathsf T}HI/2$, where
$H$ is symmetric positive semidefinite and includes the declared real
costs and net cross-firm effects. This quadratic is an additional
restriction; beliefs are not themselves social benefits. For a common
information treatment, let $I^0$ be baseline investment and
$d=\theta\pi/(1-\gamma)$. Direct expansion gives
\[
 \Delta\mathcal W
 =d(v-HI^0)^{\mathsf T}\mathbf1
 -\frac{d^2}{2}\mathbf1^{\mathsf T}H\mathbf1-c_{\mathrm{info}}, \tag{7}
\]
where $c_{\mathrm{info}}$ is the real information-production cost,
counted once. For example, take two firms, each linked to the other,
$\gamma=1/2$, $\theta=1$, $\pi=1/2$, and $a=b=0$.
Investment rises from zero to one at each firm. If validated gross value
is $Y_i=I_i$ and real research cost is $2I_i^2$, gross value rises by
two but net welfare falls by two even with zero information cost.
This uses $v=\mathbf1$ and $H=4I_n$ in (7), and all beliefs remain
between zero and one. A positive adoption or innovation-value response
therefore does not establish a positive net social gain.

\section*{Remaining gap}
The exact network formulas require known links, stable linear responses,
isolated markets, and the specified information exclusion. They neither
establish those assumptions for AI research nor identify duplication,
rival spillovers, or resource costs from adoption intentions alone.
Market-level randomization and validated value and cost measurement can
identify regime welfare without pretending that the individual Wald
ratio is a universal structural parameter. Extending identification to
unknown networks, heterogeneous nonlinear responses, and endogenous
competitor entry remains unresolved here.

\section{Completion Estimate}
49\%

This is a post hoc, heuristic estimate for the full catalogue target.
The attempt solves a linear network investment equilibrium, identifies feedback-distorted individual and market Wald ratios, and derives a welfare comparison separating investment from validated value. The full adoption-race target still needs measured belief changes and innovation outcomes, credible network and exclusion restrictions, saturation evidence, and real duplication and spillover costs.

\begin{thebibliography}{9}
\bibitem{cullen} Z. B. Cullen, E. Faia, E. Guglielminetti,
R. Perez-Truglia, and C. Rondinelli (2025), \emph{The Innovation Race:
Experimental Evidence on Advanced Technologies}, NBER Working Paper
34532, December 2025, subsequently revised. Primary abstract and record:
\url{https://www.nber.org/papers/w34532}.
\end{thebibliography}

\end{document}
